\documentclass[12pt,a4paper]{article}

\usepackage[utf8]{inputenc}
\usepackage[T2A]{fontenc}
\usepackage[russian]{babel}
\usepackage{amsmath, amssymb, amsthm}
\usepackage{geometry}
\usepackage{graphicx}
\usepackage{algorithm}
\usepackage{algorithmicx}
\usepackage{algpseudocode}
\usepackage{hyperref}

\geometry{top=2cm, bottom=2cm, left=2.5cm, right=1.5cm}

\newtheorem{theorem}{Теорема}
\newtheorem{definition}{Определение}

\title{Билет №84 \\ \small Понятие о многомерной информации. Максимальное количество измерений, в существующем физическом пространстве. Определение многомерного сигнала. (СпБПУ)}
\author{Сериков Владислав Александрович}
\date{16 мая 2026}

\begin{document}

\maketitle
\tableofcontents
\newpage

\section{Многомерная информация и многомерные сигналы}

\subsection{Понятие о многомерной информации}

\begin{definition}
\textbf{Многомерная информация} — это совокупность данных, значение которых зависит от двух и более независимых переменных (измерений, координат).
\end{definition}

В отличие от одномерной информации, которую можно представить в виде вектора или функции одной переменной (например, зависимость амплитуды от времени), многомерная информация описывается многомерными массивами (тензорами) или функциями нескольких переменных. Измерениями могут выступать пространственные координаты, время, частота, длина волны и другие физические или абстрактные параметры.

\textbf{Примеры многомерной информации:}
\begin{itemize}
    \item 2D (двумерная): цифровое изображение (зависимость интенсивности пикселя от двух пространственных координат $x$ и $y$).
    \item 3D (трехмерная): видеосигнал (две пространственные координаты и время $t$); данные магнитно-резонансной томографии (три пространственные координаты $x, y, z$).
    \item 4D (четырехмерная): объемное видео (изменение 3D-модели во времени).
    \item ND ($N$-мерная): гиперспектральные снимки (пространственные координаты, время, набор спектральных каналов).
\end{itemize}

\subsection{Максимальное количество измерений в существующем физическом пространстве}

Понятие «существующего физического пространства» зависит от рассматриваемой физической модели, описывающей Вселенную.

\begin{enumerate}
    \item \textbf{Макроскопический (наблюдаемый) уровень:}
    Классическая механика оперирует \textbf{тремя} пространственными измерениями ($x, y, z$) и одним независимым временным. Специальная и общая теории относительности (СТО и ОТО) объединяют их в единое \textbf{четырехмерное} пространство-время (пространство Минковского). На макроскопическом уровне наблюдаемыми являются ровно 4 измерения.

    \item \textbf{Микроскопический уровень (современная теоретическая физика):}
    Для объединения гравитации с квантовой механикой разработана Теория струн. Математический аппарат этой теории оказывается непротиворечивым (в частности, устраняются квантовые аномалии) только в пространствах высшей размерности:
    \begin{itemize}
        \item В теории суперструн физическое пространство содержит \textbf{10 измерений} (9 пространственных и 1 временное).
        \item В М-теории (которая обобщает 5 различных теорий суперструн) предполагается \textbf{11 измерений} (10 пространственных и 1 временное).
        \item В бозонной теории струн рассматривается \textbf{26 измерений}.
    \end{itemize}
\end{enumerate}

\textbf{Вывод:} В рамках современной фундаментальной физики максимальным количеством измерений в существующем физическом пространстве (согласно наиболее проработанной парадигме объединения взаимодействий — М-теории) считается \textbf{11 измерений}. Дополнительные 7 пространственных измерений «компактифицированы» (свернуты до планковских масштабов, например, в пространства Калаби-Яу) и не наблюдаются в макромире.

\subsection{Определение многомерного сигнала}

\begin{definition}
\textbf{Многомерный сигнал} — это физический процесс или математическая функция, несущая информацию о состоянии или поведении физической системы, которая является функцией от $M$ независимых переменных ($M \ge 2$).
\end{definition}

Математически многомерный непрерывный сигнал представляется функцией:
$$ s(\mathbf{x}) = s(x_1, x_2, \dots, x_M), \quad \mathbf{x} \in \mathbb{R}^M $$
Дискретный многомерный сигнал представляется многомерной последовательностью (решеткой):
$$ s[\mathbf{n}] = s[n_1, n_2, \dots, n_M], \quad \mathbf{n} \in \mathbb{Z}^M $$

\subsection{Связанные теоремы}

Для связи непрерывных и дискретных многомерных сигналов используется обобщение теоремы Котельникова (Найквиста-Шеннона).

\begin{theorem}[\textbf{Многомерная теорема Котельникова}]
Пусть $s(\mathbf{x})$ — непрерывный сигнал $M$ переменных, финитный в частотной области, то есть его $M$-мерное преобразование Фурье $S(\boldsymbol{\omega})$ равно нулю вне гиперпрямоугольника:
$$ S(\boldsymbol{\omega}) = 0 \quad \text{для} \quad |\omega_k| > \Omega_k, \quad k = 1, 2, \dots, M, $$
где $\boldsymbol{\omega} = (\omega_1, \dots, \omega_M)$ — вектор пространственных частот.
Тогда сигнал $s(\mathbf{x})$ может быть полностью и без искажений восстановлен по своим дискретным отсчетам $s(\mathbf{n}\Delta\mathbf{x}) = s(n_1 \Delta x_1, \dots, n_M \Delta x_M)$, взятым с шагами:
$$ \Delta x_k \le \frac{\pi}{\Omega_k}, \quad k = 1, \dots, M. $$
\end{theorem}

\begin{proof}
Представим $M$-мерное преобразование Фурье сигнала:
$$ S(\boldsymbol{\omega}) = \int_{\mathbb{R}^M} s(\mathbf{x}) e^{-j \boldsymbol{\omega}^T \mathbf{x}} d\mathbf{x} $$
Так как спектр ограничен областью $B = [-\Omega_1, \Omega_1] \times \dots \times [-\Omega_M, \Omega_M]$, мы можем периодически продолжить его на все пространство частот с периодами $2\Omega_k$. Полученный периодический спектр $\tilde{S}(\boldsymbol{\omega})$ можно разложить в многомерный ряд Фурье:
$$ \tilde{S}(\boldsymbol{\omega}) = \sum_{n_1=-\infty}^{\infty} \dots \sum_{n_M=-\infty}^{\infty} C_{\mathbf{n}} e^{-j \sum_{k=1}^M \frac{\pi}{\Omega_k} n_k \omega_k} $$
Найдем коэффициенты ряда Фурье $C_{\mathbf{n}}$:
$$ C_{\mathbf{n}} = \frac{1}{\prod_{k=1}^M 2\Omega_k} \int_{-\Omega_1}^{\Omega_1} \dots \int_{-\Omega_M}^{\Omega_M} \tilde{S}(\boldsymbol{\omega}) e^{j \sum_{k=1}^M \frac{\pi}{\Omega_k} n_k \omega_k} d\boldsymbol{\omega} $$
Заметим, что обратное преобразование Фурье исходного сигнала $s(\mathbf{x})$ в точке $\mathbf{x} = \left(\frac{n_1 \pi}{\Omega_1}, \dots, \frac{n_M \pi}{\Omega_M}\right)$ в точности совпадает (с точностью до постоянного множителя) с выражением для интеграла:
$$ s\left(\frac{n_1 \pi}{\Omega_1}, \dots, \frac{n_M \pi}{\Omega_M}\right) = \frac{1}{(2\pi)^M} \int_{-\Omega_1}^{\Omega_1} \dots \int_{-\Omega_M}^{\Omega_M} S(\boldsymbol{\omega}) e^{j \sum_{k=1}^M \frac{\pi}{\Omega_k} n_k \omega_k} d\boldsymbol{\omega} $$
Отсюда следует:
$$ C_{\mathbf{n}} = \frac{(2\pi)^M}{\prod_{k=1}^M 2\Omega_k} s\left(\frac{n_1 \pi}{\Omega_1}, \dots, \frac{n_M \pi}{\Omega_M}\right) $$
Таким образом, спектр $S(\boldsymbol{\omega})$ на базовом интервале $B$ полностью определяется отсчетами функции $s(\mathbf{x})$. Подставляя ряд Фурье в формулу обратного преобразования Фурье и интегрируя по прямоугольной области, получаем формулу восстановления (интерполяционный ряд):
$$ s(\mathbf{x}) = \sum_{n_1=-\infty}^{\infty} \dots \sum_{n_M=-\infty}^{\infty} s\left(\frac{n_1 \pi}{\Omega_1}, \dots, \frac{n_M \pi}{\Omega_M}\right) \prod_{k=1}^M \text{sinc}\left(\frac{\Omega_k x_k - n_k \pi}{\pi}\right), $$
где $\text{sinc}(u) = \frac{\sin(\pi u)}{\pi u}$.
Доказательство завершено: непрерывный многомерный сигнал восстановлен из дискретных отсчетов.
\end{proof}

\subsection{Алгоритмы обработки: Многомерное Быстрое Преобразование Фурье}

Обработка многомерных сигналов требует эффективных алгоритмов. Наиболее важным алгоритмом является $M$-мерное дискретное преобразование Фурье (ДПФ). Благодаря свойству сепарабельности базисных функций, многомерное преобразование сводится к последовательности одномерных преобразований.

Ниже приведен алгоритм для двумерного случая (2D БПФ), применяемый к матрице изображения.

\begin{algorithm}[H]
\caption{Алгоритм 2D БПФ (Row-Column Decomposition)}
\begin{algorithmic}[1]
\Require Матрица $X$ размером $N \times M$ (двумерный дискретный сигнал).
\Ensure Матрица $Y$ размером $N \times M$ (пространственный спектр).
\State Создать промежуточную матрицу $X'$ размером $N \times M$.
\For{$i = 0$ \textbf{to} $N-1$} \Comment{Обработка строк}
    \State Извлечь $i$-ю строку: $\mathbf{r}_i = X[i, :]$
    \State Применить одномерное БПФ: $\mathbf{r}'_i = \text{1D\_FFT}(\mathbf{r}_i)$
    \State Сохранить результат: $X'[i, :] = \mathbf{r}'_i$
\EndFor
\For{$j = 0$ \textbf{to} $M-1$} \Comment{Обработка столбцов промежуточной матрицы}
    \State Извлечь $j$-й столбец: $\mathbf{c}_j = X'[:, j]$
    \State Применить одномерное БПФ: $\mathbf{c}'_j = \text{1D\_FFT}(\mathbf{c}_j)$
    \State Сохранить результат: $Y[:, j] = \mathbf{c}'_j$
\EndFor
\State \Return $Y$
\end{algorithmic}
\end{algorithm}

\textbf{Минимальный пример выполнения алгоритма:}
Рассмотрим сигнал размером $2 \times 2$:
$$ X = \begin{pmatrix} 1 & 2 \\ 3 & 4 \end{pmatrix} $$
Используем 2-точечное одномерное БПФ, где $\text{1D\_FFT}([a, b]) = [a+b, a-b]$.

\textit{Шаг 1. Обработка строк:}
\begin{itemize}
    \item Строка 0: $[1, 2] \rightarrow [1+2, 1-2] = [3, -1]$
    \item Строка 1: $[3, 4] \rightarrow [3+4, 3-4] = [7, -1]$
\end{itemize}
Промежуточная матрица:
$$ X' = \begin{pmatrix} 3 & -1 \\ 7 & -1 \end{pmatrix} $$

\textit{Шаг 2. Обработка столбцов (матрицы $X'$):}
\begin{itemize}
    \item Столбец 0: $[3, 7]^T \rightarrow [3+7, 3-7]^T = [10, -4]^T$
    \item Столбец 1: $[-1, -1]^T \rightarrow [-1+(-1), -1-(-1)]^T = [-2, 0]^T$
\end{itemize}
Итоговый спектр:
$$ Y = \begin{pmatrix} 10 & -2 \\ -4 & 0 \end{pmatrix} $$

Этот алгоритм имеет сложность $O(N M \log(NM))$, что существенно быстрее прямого вычисления двумерного ДПФ по определению со сложностью $O(N^2 M^2)$.

\end{document}
